\chapter{Методы используемые в исследовании спектров молекул}\label{ch:methods}

\section{Оператор Гамильтона молекулы}\label{sec:methods/hamilton}

\subsection{Роль гамильтониана в квантово-механическом описании молекул}

Квантово-механическое описание молекулярных систем основано на использовании
оператора Гамильтона, определяющего полную энергию системы.
В стационарной квантовой механике состояние молекулы описывается
волновой функцией $\Psi$, удовлетворяющей уравнению Шрёдингера

\begin{equation}
\hat{H}\Psi = E\Psi ,
\end{equation}

где $\hat{H}$ — оператор Гамильтона, а $E$ — собственное значение энергии.

Собственные значения гамильтониана соответствуют энергетическим уровням
молекулы, а переходы между ними приводят к появлению линий в спектрах
поглощения и излучения. Таким образом, анализ молекулярных спектров
фактически сводится к исследованию структуры собственных значений
гамильтониана.

Формирование гамильтонова подхода в молекулярной физике связано
с развитием квантовой механики и спектроскопии в первой половине XX века.
Фундаментальные основы теории колебательно-вращательных спектров
были сформулированы в классических работах
\cite{nielsen,Amat1971,Makushkin1984}.
Эти исследования показали, что даже относительно простые молекулярные
системы обладают сложной энергетической структурой,
которая определяется взаимосвязью вращательного и колебательного движения.

\subsection{Полный гамильтониан молекулы}

В нерелятивистском приближении гамильтониан молекулы,
состоящей из $N$ электронов и $M$ атомных ядер, можно записать как

\begin{equation}
\hat{H} =
-\sum_{i=1}^{N}\frac{\hbar^2}{2m_e}\nabla_i^2
-\sum_{A=1}^{M}\frac{\hbar^2}{2M_A}\nabla_A^2
+\sum_{i<j}\frac{e^2}{r_{ij}}
-\sum_{i,A}\frac{Z_A e^2}{r_{iA}}
+\sum_{A<B}\frac{Z_A Z_B e^2}{R_{AB}} .
\end{equation}

Здесь первые два члена описывают кинетическую энергию электронов
и атомных ядер, а остальные члены учитывают кулоновские взаимодействия
между всеми частицами системы.

Прямое решение уравнения Шрёдингера с таким гамильтонианом
для многоатомных молекул представляет собой чрезвычайно сложную задачу.
Поэтому в практических расчётах используются приближённые методы,
основой которых является приближение Борна–Оппенгеймера.
В этом приближении движение электронов и ядер рассматривается раздельно,
что позволяет сначала решить электронную задачу,
а затем исследовать движение ядер в эффективном потенциале.

После отделения движения центра масс остаётся задача описания
внутреннего движения ядер, включающего вращательные
и колебательные степени свободы.

\subsection{Колебательно-вращательный гамильтониан}

Внутреннее движение ядер молекулы характеризуется
$3N-6$ колебательными координатами
(для нелинейной молекулы) и тремя вращательными степенями свободы.
Соответствующий гамильтониан можно представить в виде

\begin{equation}
\hat{H} =
\hat{H}_{vib} +
\hat{H}_{rot} +
\hat{H}_{vr},
\end{equation}

где $\hat{H}_{vib}$ описывает колебательное движение,
$\hat{H}_{rot}$ — вращательное движение,
а $\hat{H}_{vr}$ учитывает взаимодействие между ними.

В гармоническом приближении колебательная часть имеет вид

\begin{equation}
\hat{H}_{vib} =
\frac{1}{2}\sum_k
\left(P_k^2 + \omega_k^2 Q_k^2 \right),
\end{equation}

где $Q_k$ — нормальные координаты,
$P_k$ — сопряжённые импульсы,
а $\omega_k$ — частоты нормальных колебаний.

Вращательная часть гамильтониана для нелинейной молекулы
может быть записана как

\begin{equation}
\hat{H}_{rot} =
A\hat{J}_a^2 +
B\hat{J}_b^2 +
C\hat{J}_c^2 ,
\end{equation}

где $A$, $B$ и $C$ — вращательные постоянные,
связанные с главными моментами инерции молекулы.

\subsection{Операторы углового момента}

Вращательное движение молекулы описывается оператором
полного углового момента $\hat{\mathbf{J}}$.
Его компоненты $\hat{J}_a$, $\hat{J}_b$, $\hat{J}_c$
в молекулярной системе координат удовлетворяют
коммутационным соотношениям

\begin{equation}
[\hat{J}_i,\hat{J}_j] = i\hbar \varepsilon_{ijk}\hat{J}_k .
\end{equation}

Собственные функции операторов $\hat{J}^2$ и $\hat{J}_z$
определяются квантовыми числами $J$ и $M$:

\begin{align}
\hat{J}^2 |J,M\rangle &= \hbar^2 J(J+1)|J,M\rangle, \\
\hat{J}_z |J,M\rangle &= \hbar M |J,M\rangle .
\end{align}

Квантовое число $J$ характеризует полный момент вращения молекулы,
а $M$ определяет его проекцию на выбранную ось.

Для симметричных молекул дополнительно вводится
квантовое число $K$, определяющее проекцию
углового момента на ось симметрии молекулы.

\subsection{Симметрия волновых функций}

Волновая функция молекулы может быть представлена
в виде произведения

\begin{equation}
\Psi = \Psi_{el}\Psi_{vib}\Psi_{rot}.
\end{equation}

Симметрия каждой из этих функций определяется
точечной группой симметрии молекулы.
Использование свойств симметрии позволяет существенно
уменьшить размерность квантово-механической задачи,
поскольку гамильтониан коммутирует с операциями
группы симметрии.

Симметрия состояний играет важную роль
в формулировке правил отбора для спектроскопических переходов
и при классификации энергетических уровней молекулы.

\subsection{Вывод гамильтониана Уотсона}

При описании движения ядер в молекуле
необходимо учитывать как вращение,
так и колебания атомов.
Полное разделение этих движений невозможно,
однако их связь можно минимизировать
путём выбора специальной системы координат,
известной как система Экарта \cite{Eckart1935,BunkerJensen1998}.
В такой системе координат оси связаны с молекулой
и изменяются вместе с её колебаниями.

После применения условий Экарта
гамильтониан движения ядер можно преобразовать
в более удобную форму.
В 1968 году Дж.~Уотсон предложил преобразование,
позволяющее существенно упростить
колебательно-вращательный гамильтониан молекулы
\cite{Watson1968}.

В результате гамильтониан движения ядер
можно записать в виде

\begin{equation}
\hat{H} =
-\frac{\hbar^2}{2}\sum_{s=1}^{3N-6}
\frac{\partial^2}{\partial q_s^2}
+
\frac{1}{2}
\sum_{\alpha\beta}
\mu_{\alpha\beta}
(\hat{J}_\alpha - \hat{\pi}_\alpha)
(\hat{J}_\beta - \hat{\pi}_\beta)
+V(q_1,\ldots,q_{3N-6}),
\end{equation}

где

\begin{itemize}
\item $q_s$ — нормальные колебательные координаты,
\item $\hat{J}_\alpha$ — компоненты оператора углового момента,
\item $\hat{\pi}_\alpha$ — операторы кориолисового взаимодействия,
\item $\mu_{\alpha\beta}$ — элементы обратного тензора инерции,
\item $V$ — потенциальная энергия.
\end{itemize}

Полученный гамильтониан известен как гамильтониан Уотсона
и широко используется при анализе
колебательно-вращательных спектров молекул
\cite{Watson1968,Watson1977}.

\subsection{Переход к матричному представлению гамильтониана}

Для практического вычисления энергетических уровней
гамильтониан представляется в виде матрицы
в выбранном базисе колебательно-вращательных функций

\begin{equation}
H_{ij} =
\langle \Phi_i | \hat{H} | \Phi_j \rangle .
\end{equation}

Диагонализация этой матрицы позволяет получить
собственные значения энергии и соответствующие
волновые функции молекулы.

Построение эффективных базисных функций
и вычисление матричных элементов гамильтониана
рассматриваются в следующем разделе,
где будет подробно описано
матричное представление гамильтониана
и методы его использования
при анализе колебательно-вращательных спектров.

\section{Приближение Борна-Оппенгеймера}

\subsection{Физический смысл и предпосылки приближения}

При квантово-механическом описании молекулы необходимо одновременно
учитывать движение электронов и атомных ядер. Полная волновая функция
молекулярной системы зависит от координат всех частиц, а соответствующее
уравнение Шрёдингера содержит как электронную, так и ядерную кинетическую
энергию, а также все кулоновские взаимодействия. Для многоатомных
молекул такая задача в исходной постановке обладает чрезвычайно высокой
размерностью, вследствие чего её прямое решение оказывается практически
невыполнимым даже для сравнительно простых систем.

Ключевое упрощение связано с тем, что атомные ядра существенно тяжелее
электронов. Именно это различие масштабов масс приводит к различию
характерных времён движения соответствующих подсистем. Электронное
облако успевает перестроиться значительно быстрее, чем заметно изменяется
положение ядер. Иными словами, для электронной подсистемы конфигурацию
ядер в первом приближении можно считать фиксированной, тогда как для
ядерного движения электронная энергия играет роль эффективного потенциала.
Такой подход и составляет физическую основу приближения
Борна--Оппенгеймера.

Исторически данное приближение стало одним из важнейших результатов
раннего этапа развития квантовой механики. Оно позволило перейти
от формально точного, но практически непригодного полного уравнения
Шрёдингера к иерархической схеме описания молекулы, в которой
электронная, колебательная и вращательная задачи рассматриваются
последовательно. Благодаря этому приближение Борна--Оппенгеймера
фактически стало фундаментом современной квантовой химии,
молекулярной физики и спектроскопии высокого разрешения.

\subsection{Разделение полной молекулярной задачи}

Пусть $\mathbf{r}$ обозначает совокупность электронных координат,
а $\mathbf{R}$ --- совокупность координат атомных ядер.
Тогда полная волновая функция молекулы имеет вид
$\Psi(\mathbf{r},\mathbf{R})$ и удовлетворяет стационарному уравнению

\begin{equation}
\hat{H}\Psi(\mathbf{r},\mathbf{R})=E\Psi(\mathbf{r},\mathbf{R}).
\end{equation}

В рамках приближения Борна--Оппенгеймера предполагается, что полную
волновую функцию можно представить в виде произведения электронной
и ядерной частей:

\begin{equation}
\Psi(\mathbf{r},\mathbf{R})
=
\psi_{el}(\mathbf{r};\mathbf{R})\chi(\mathbf{R}).
\end{equation}

Здесь электронная функция $\psi_{el}(\mathbf{r};\mathbf{R})$
зависит от координат ядер параметрически:
при решении электронной задачи положения ядер фиксируются,
а их координаты входят в уравнение лишь как внешние параметры.
Ядерная функция $\chi(\mathbf{R})$ описывает уже более медленное
движение самих ядер.

Подстановка такого разложения в полное уравнение Шрёдингера приводит
к естественному разделению задачи на два этапа.
Сначала решается электронное уравнение при фиксированной конфигурации ядер:

\begin{equation}
\hat{H}_{el}(\mathbf{r};\mathbf{R})
\psi_{el}(\mathbf{r};\mathbf{R})
=
E_{el}(\mathbf{R})\psi_{el}(\mathbf{r};\mathbf{R}),
\end{equation}

где $\hat{H}_{el}$ включает кинетическую энергию электронов,
электрон-электронное взаимодействие, электрон-ядерное взаимодействие
и межъядерное отталкивание как параметрическую величину.
Для каждой фиксированной конфигурации ядер получается собственное
значение $E_{el}(\mathbf{R})$, которое зависит от геометрии молекулы.

Далее найденная электронная энергия используется в уравнении
для ядерного движения:

\begin{equation}
\left[
-\sum_A\frac{\hbar^2}{2M_A}\nabla_A^2
+E_{el}(\mathbf{R})
\right]\chi(\mathbf{R})
=
E\chi(\mathbf{R}),
\end{equation}

если не учитывать более тонкие поправки, возникающие из-за
неполного разделения электронного и ядерного движения.
Тем самым исходная многочастичная задача заменяется последовательным
решением электронной и ядерной задач, что делает возможным
практический анализ молекулярных состояний.

\subsection{Потенциальные поверхности энергии и их роль}

Функция $E_{el}(\mathbf{R})$, получаемая как результат решения
электронного уравнения, называется потенциальной поверхностью энергии.
Именно она определяет, в каком эффективном поле движутся ядра,
а значит, задаёт геометрию молекулы, её колебательные свойства
и значительную часть наблюдаемой спектроскопической структуры.

Минимум потенциальной поверхности соответствует равновесной конфигурации
молекулы. В окрестности этого минимума движение ядер может быть
рассмотрено как малые отклонения от равновесия, что приводит
к введению нормальных координат и гармоническому приближению
для колебаний. Кривизна поверхности вблизи минимума определяет
гармонические частоты, а члены более высокого порядка описывают
ангармоничность и взаимодействие колебательных мод.

Если в системе существует несколько электронных состояний,
то каждому из них соответствует своя потенциальная поверхность.
В большинстве спектроскопических задач высокого разрешения
рассматривается движение ядер на одной выбранной электронной поверхности,
обычно соответствующей основному электронному состоянию.
Именно в этом случае приближение Борна--Оппенгеймера проявляет
свою наибольшую эффективность.

Следует подчеркнуть, что в молекулярной спектроскопии
потенциальная поверхность энергии является не просто промежуточной
математической конструкцией, а фактически центральным объектом
описания внутренней динамики молекулы. Через неё определяются
равновесные структурные параметры, колебательные уровни,
резонансные взаимодействия и, в конечном счёте, положение
наблюдаемых спектральных линий.

\subsection{Адиабатическое приближение и структура молекулярной энергии}

При использовании приближения Борна--Оппенгеймера молекулярную энергию
естественно рассматривать как величину, распадающуюся на несколько
компонент, соответствующих различным типам движения. В первом
приближении полная энергия молекулы может быть представлена как сумма
электронного, колебательного и вращательного вкладов. Такое разложение
не является строго точным, однако именно оно лежит в основе всей
классификации молекулярных уровней и спектров.

Сначала формируется электронная структура молекулы, определяемая
решением электронной задачи. Затем на фиксированной электронной
поверхности возникают колебательные уровни, связанные с движением
ядер относительно равновесной конфигурации. Наконец, каждый
колебательный уровень расщепляется на вращательные подуровни.
Именно эта иерархия масштабов энергии делает возможным
последовательное построение молекулярного гамильтониана
и поэтапный анализ спектров.

Такой подход особенно важен для колебательно-вращательной
спектроскопии. Наблюдаемые в инфракрасных спектрах переходы,
как правило, происходят в пределах одного электронного состояния
и отражают изменение колебательных и вращательных квантовых чисел.
Следовательно, сам факт возможности анализировать эти переходы
независимо от электронной структуры в значительной степени основан
именно на приближении Борна--Оппенгеймера.

\subsection{Математические основания и малый параметр}

С формальной точки зрения приближение Борна--Оппенгеймера опирается
на существование малого параметра, связанного с отношением массы
электрона к характерной массе ядра. Благодаря этому полное уравнение
Шрёдингера может быть исследовано асимптотически, а электронное
и ядерное движения оказываются разделимыми в нулевом порядке
по указанному малому параметру.

Именно в этом смысле приближение Борна--Оппенгеймера не является
лишь качественным рассуждением о ``быстрых'' электронах и
``медленных'' ядрах, а представляет собой вполне определённую
схему последовательного приближения. В первом порядке получаются
разделённые электронная и ядерная задачи. В более высоких порядках
возникают поправки, связанные с тем, что электронная волновая функция
всё же зависит от координат ядер не только параметрически, но и
через производные по этим координатам.

Эти дополнительные члены обычно малы, однако в ряде случаев
они могут оказывать заметное влияние на спектр. Поэтому
в прикладной спектроскопии важно не только использовать
приближение Борна--Оппенгеймера как исходный уровень описания,
но и понимать границы его применимости.

\subsection{Нарушение приближения и неадиабатические эффекты}

Хотя приближение Борна--Оппенгеймера чрезвычайно эффективно,
оно не является строго точным. Основная причина заключается в том,
что электронная функция зависит от координат ядер, а потому
при действии ядерного кинетического оператора появляются дополнительные
члены, связывающие разные электронные состояния. Эти члены и образуют
так называемые неадиабатические связи.

В обычных условиях, когда электронные состояния хорошо разделены по энергии,
такие связи малы, и ими можно пренебречь. Однако ситуация меняется,
если потенциальные поверхности сближаются или пересекаются.
В этом случае разделение электронного и ядерного движения ухудшается,
а простая одноповерхностная картина становится недостаточной.
Подобные эффекты особенно важны в фотохимии, в теории элементарных
реакций, а также при анализе некоторых тонких спектроскопических
особенностей высоковозбуждённых состояний.

Для задач, связанных с высокоразрешённой инфракрасной спектроскопией
устойчивых молекул в основном электронном состоянии, нарушение
приближения Борна--Оппенгеймера обычно не является определяющим.
Тем не менее именно понимание существования таких поправок
объясняет, почему даже очень точные эффективные гамильтонианы
следует рассматривать как приближённые модели, а не как абсолютно
полное описание молекулярной динамики.

\subsection{Значение приближения Борна--Оппенгеймера для молекулярной спектроскопии}

Для молекулярной спектроскопии приближение Борна--Оппенгеймера
имеет принципиальное значение. Именно оно позволяет перейти
от общей многочастичной задачи к последовательному описанию
электронных, колебательных и вращательных состояний, что делает
возможной интерпретацию экспериментальных спектров в терминах
квантовых чисел и эффективных спектроскопических параметров.

В частности, после фиксации электронной поверхности энергии
задача сводится к исследованию движения ядер. Далее вводятся
нормальные координаты, строится колебательный гамильтониан,
а затем учитывается вращательное движение и его связь с колебаниями.
Именно в этой логике формируются как строгие квантово-механические,
так и полуэмпирические методы анализа спектров высокого разрешения.

Таким образом, приближение Борна--Оппенгеймера выступает не просто
как техническое упрощение исходного уравнения Шрёдингера, а как
фундаментальный принцип организации молекулярной теории.
Оно связывает электронную структуру молекулы с её ядерной динамикой
и тем самым создаёт теоретическую основу для последующего рассмотрения
колебательных состояний, колебательно-вращательного взаимодействия
и построения эффективных гамильтонианов, используемых
в практической спектроскопии.

В дальнейшем, опираясь на данное приближение, можно перейти
к описанию колебательного движения ядер в окрестности равновесной
конфигурации молекулы. Именно этот этап приводит к введению
нормальных координат, колебательных волновых функций и,
в конечном счёте, к матричному представлению молекулярного
гамильтониана, необходимому для анализа экспериментальных спектров.

\section{Теория изотопозамещения}

\section{Неприводимые тензорные операторы}

\section{Методы восстановления потенциальной функции молекулы на основе экспериментальных данных}

\subsection{Методы ab initio}

\subsection{Полуэмпирические методы}

